% =====================================================================
%  30.5  The method boundary: labelled vs unlabelled, and where each
%  tool of this chapter stops applying.
%  Exports: sec:c32-boundary  def:c32-orbit  thm:c32-orbit  def:c32-pick
%  fig:c32-orbit
%  WORKED examples only.  Hands off to Ch.28 (shapes), Ch.26 (virtual
%  trees), Ch.31 (MST).  The reader leaves Ch.30 with the split "count
%  sequences / take a determinant / split into indicators / hand off".
% =====================================================================
\section{The Method Boundary: Labelled vs.\ Unlabelled}\label{sec:c32-boundary}

\begin{previouslybox}[Ch.28 compared shapes; Ch.30 counted labelings]
Chapter~28 (\chref{c22}) took two trees and decided whether they were the
\emph{same shape}, running AHU or a tree hash to canonicalise.  This chapter
counts \emph{labelled} trees, where two assignments of the same shape to the
vertex set \([n]\) are different objects.  The two views are reconciled by a
single group-action identity: the size of the orbit of a labelled tree under
\(S_n\) is \(n!/|\mathrm{aut}(T)|\).  Section \(\secref{c32-boundary}\) makes that
precise, hands the unlabelled case back to Ch.~28, and lists three nearby
topics -- virtual trees (Ch.~26 \chref{c30}) and MST (Ch.~31 \chref{c17})
among them -- that are \emph{not} this chapter's territory.
\end{previouslybox}

\begin{topicintro}
Two questions that sound identical and are not: "How many trees are there on
\(n\) vertices?" and "How many trees are there with \(n\) vertices?".  The
first is Cayley's \(n^{n-2}\) (labelled: which subset of \([n]\times[n]\) is
the edge set).  The second is the count of \emph{unlabelled tree shapes}
(A000055 in OEIS); the answer is \(1, 1, 1, 2, 3, 6, 11, 23, 47, 106,\dots\)
and has no closed form -- it is computed by P\'olya enumeration or by the
rooted-shape recursion of Ch.~28.  Deciding which question a statement asks
is the single most important reading skill of this section.

The bridge between the two is the orbit-stabiliser theorem
(\(\thmref{c32-orbit}\)): \(S_n\) acts on \(\mathcal{T}_n\) by relabelling,
and each unlabelled shape \(S\) contributes \(n!/|\mathrm{aut}(S)|\)
labelled trees to the total.  Cayley's number is therefore \(\sum_{S\text{
shape on }n} n!/|\mathrm{aut}(S)|\).  This identity is why the unlabelled
enumeration is harder -- \(|\mathrm{aut}(S)|\) depends on the shape's
symmetry structure, which Ch.~28 is precisely the machinery to compute.

The rest of this section is a \emph{pick list}: four questions you will meet
that are \emph{not} Chapter~30's to answer, with a one-line pointer to the
chapter that owns each.  The chapter ends with the Summit CF 917D
(\(\probref{c32-summit}\)), which is a genuine Ch.~30 question mark,
unlabelled on the surface but labelled underneath.

\textbf{Prerequisites.} Ch.~28 (\chref{c22}) for \(|\mathrm{aut}(T)|\) via
AHU; Ch.~26 (\chref{c30}) for virtual trees; Ch.~31 (\chref{c17}) for MST.

\textbf{After this section you can}
\begin{itemize}[nosep]
  \item distinguish labelled from unlabelled counting in a problem statement; \hfill\lv{4}
  \item apply orbit-stabiliser to convert between the two views; \hfill\lv{5}
  \item name the three nearest topics this chapter does \emph{not} own and cite their chapters. \hfill\lv{4}
\end{itemize}
\end{topicintro}

% ---------------------------------------------------------------
\subsection{The labelled vs.\ unlabelled split}
% ---------------------------------------------------------------

\begin{definition}{Automorphism of a tree}{c32-orbit}
An \term{automorphism} of a labelled tree \(T\in\mathcal{T}_n\) is a
permutation \(\sigma\in S_n\) with \(\{u,v\}\in E(T)\iff \{\sigma(u),
\sigma(v)\}\in E(T)\).  The automorphism group \(\mathrm{aut}(T)\) is the
set of such \(\sigma\); its size measures the shape's internal symmetry.
\end{definition}

\begin{theorem}{Labelled count via unlabelled shapes}{c32-orbit}
Let \(\mathcal{S}_n\) be the set of unlabelled tree shapes on \(n\) vertices.
Then \(|\mathcal{T}_n|=\sum_{S\in\mathcal{S}_n} n!/|\mathrm{aut}(S)|\).
\end{theorem}

\begin{proofbox}
The symmetric group \(S_n\) acts on \(\mathcal{T}_n\) by relabelling.  Two
labelled trees \(T,T'\) are in the same orbit iff they share an unlabelled
shape.  By orbit-stabiliser, the size of the orbit containing \(T\) is
\([S_n:\mathrm{Stab}(T)]=n!/|\mathrm{Stab}(T)|\), and \(\mathrm{Stab}(T)=
\mathrm{aut}(T)\).  Summing orbit sizes over the set of shapes gives the
total \(|\mathcal{T}_n|\). \(\Box\)
\end{proofbox}

\subsection*{Worked: reconciling \(n=4\)}

Ch.~28 has \(2\) unlabelled tree shapes on \(4\) vertices: the path \(P_4\)
and the star \(K_{1,3}\).  \(|\mathrm{aut}(P_4)|=2\) (the identity and the
end-to-end flip), \(|\mathrm{aut}(K_{1,3})|=6\) (all permutations of the
three leaves).  Orbit sizes: \(4!/2 = 12\) labelled copies of \(P_4\), and
\(4!/6 = 4\) labelled copies of \(K_{1,3}\).  Sum \(=16\), matching Cayley
\(4^{4-2}=16\).  Both views agree.

\begin{figure}[H]
\centering
\begin{tikzpicture}[diagram, scale=0.9, transform shape,
   tnode/.style={circle, draw=primaryblue!70, fill=primaryblue!8, line width=0.7pt,
                 minimum size=6mm, font=\scriptsize\bfseries, inner sep=1pt}]
  \node[tnode] (p1) at (0,0) {1};
  \node[tnode] (p2) at (1,0) {2};
  \node[tnode] (p3) at (2,0) {3};
  \node[tnode] (p4) at (3,0) {4};
  \draw (p1)--(p2)--(p3)--(p4);
  \node[annot, font=\scriptsize] at (1.5,-0.9) {\(P_4\): aut size 2, orbit 12};
  \node[tnode] (c) at (6,0) {1};
  \node[tnode] (l1) at (5,1) {2};
  \node[tnode] (l2) at (7,1) {3};
  \node[tnode] (l3) at (7,-1) {4};
  \draw (c)--(l1) (c)--(l2) (c)--(l3);
  \node[annot, font=\scriptsize] at (6,-2) {\(K_{1,3}\): aut size 6, orbit 4};
\end{tikzpicture}
\caption{The two unlabelled tree shapes on \(4\) vertices.  Their orbit
sizes add to Cayley's \(16\).}
\label{fig:c32-orbit}
\end{figure}

% ---------------------------------------------------------------
\subsection{The pick list: what this chapter does not own}
% ---------------------------------------------------------------

\begin{definition}{Method-boundary pick list}{c32-pick}
Read the statement and pick:
\begin{itemize}[nosep]
  \item \textbf{Labelled count, closed form possible} -- Pr\"ufer
        (\(\secref{c32-prufer}\)) or Cayley variants
        (\(\secref{c32-cayley}\)).
  \item \textbf{Labelled count inside a \emph{given} graph} -- Matrix-Tree
        (\(\secref{c32-matrix}\)).
  \item \textbf{Expected value of a random process on a tree} -- linearity
        of expectation (\(\secref{c32-expect}\)).
  \item \textbf{Unlabelled shape census, automorphism size, tree isomorphism}
        -- Ch.~28 (\chref{c22}); use \(\thmref{c32-orbit}\) to reconcile
        against a labelled count if the problem mixes the two views.
  \item \textbf{Minimum spanning tree of a weighted graph} -- Ch.~31
        (\chref{c17}); the "tree" is a \emph{specific} one, not a count.
  \item \textbf{Tree over a marked vertex subset, distance queries} --
        Ch.~26 (\chref{c30}) virtual trees; the object is a small tree
        built per query, not a family.
\end{itemize}
\end{definition}

\begin{insightbox}[The one question you should be able to answer]
Given a problem statement, is the answer \emph{a number mod a prime}, a
\emph{real number}, or \emph{a specific tree}?  A number mod a prime means
you are counting labelled or spanning trees (Ch.~30).  A real number means
expectation (Ch.~30.4 is your template).  A specific tree means construction,
which is Ch.~26 or Ch.~31 territory.  This three-way split is the whole of
\(\secref{c32-boundary}\); the rest of the book elaborates.
\end{insightbox}

\begin{notebox}[A cautionary "count unlabelled trees" example]
"How many unlabelled tree shapes on \(n=8\) vertices?" is a Ch.~28 question;
the answer is \(23\).  No Ch.~30 tool applies -- Pr\"ufer is a bijection on
\emph{labelled} trees and Matrix-Tree counts \emph{labelled} spanning trees
of a specific graph.  Trying to "divide Cayley's \(8^6\) by \(8!\)" is a
classic wrong move; the orbit sizes vary and you cannot factor out a
uniform \(n!\).  Ch.~28's AHU or P\'olya machinery handles the enumeration.
\end{notebox}

\begin{recallbox}[30.5 in one breath]
\emph{Labelled trees form an \(S_n\)-set; unlabelled shapes are its orbits,
and \(\thmref{c32-orbit}\) is the bridge.  Cayley counts labels, AHU counts
shapes, and mixing the two without the orbit correction is the classic
error.  Virtual trees (Ch.~26 \chref{c30}) and MST (Ch.~31 \chref{c17}) are
one step outside this chapter's territory -- leave the answer space and you
leave Ch.~30.}
\selftest{
  \item Reconcile \(|\mathcal{T}_4| = 16\) via the two shapes on \(4\)
        vertices.
  \item Explain why dividing Cayley by \(n!\) is wrong for unlabelled
        enumeration.
  \item Name one Ch.~26 problem that is \emph{not} a Ch.~30 problem.
}
\end{recallbox}
